% clase article a una cara, con letra de 12pt, en papel A4
% y con p�gina separada para el t�tulo
\documentclass[oneside,12pt,a4paper,titlepage]{article}

% Carga de paquetes necesarios. "customarticle" es un paquete personalizado
\usepackage[spanish]{babel} 
\RequirePackage[T1]{fontenc}
\RequirePackage[ansinew]{inputenx}
\usepackage{array}
\usepackage{graphicx}
\usepackage{pifont}
\usepackage[usenames,dvipsnames]{color}
\usepackage{amsmath}
\usepackage{hyperref}
\usepackage{colortbl}
\usepackage{makeidx}
\hypersetup{bookmarksopen,bookmarksopenlevel=3,linktocpage,colorlinks,urlcolor=blue,citecolor=blue,linkcolor=blue,filecolor=blue,pdfnewwindow,
	pdftitle={Cuadro de Mando para la Visualizaci�n de la Calidad del Software, Integrado con un Entorno de Medici�n}, 
	pdfsubject={Proyecto Fin de Carrera}
	pdfauthor={Jose Domingo L�pez L�pez}}
\usepackage[spanish,caption,counter,title,fancy]{customarticle}

\makeindex

% Macro para insertar una imagen
%       Uso: \imagen{nombreFichero}{Factor escala}{Caption (leyenda)}{Label (identificador para referenciarla)}
\def\imagen#1#2#3#4{
 \begin{figure}[h]
 \begin{center}
   \includegraphics[keepaspectratio,scale=#2]{#1}
 \caption {#3}
 \label{#4}
 \end{center}
 \end{figure}
}

% Definici�n de una lista personalizada 
\newenvironment{myitemize}%
{\begin{list}{\textbullet}%
{\settowidth{\labelwidth}{\textbullet}
\addtolength{\topsep}{-1.5ex}
\addtolength{\itemsep}{\dimexpr 0ex plus 0.25ex minus 0.25ex}
\setlength{\partopsep}{0.5ex}
}}%
{\end{list}}

\begin{document}

% En las p�ginas de portada e �ndices, no hay encabezado ni pie de p�gina
\pagestyle{empty}
\begin{titlepage}
	\begin{center}
		\includegraphics[scale=0.25, keepaspectratio]{./images/esi_bw.png} \\
		\vspace{20mm}
  	{\large \textbf{UNIVERSIDAD DE CASTILLA-LA MANCHA}} \\
  	\bigskip
  	{\large \textbf{ESCUELA SUPERIOR DE INFORM�TICA}} \\
  	\vspace{25mm}
  	{\large \textbf{INGENIER�A}}\\
  	\bigskip
  	{\large \textbf{EN INFORM�TICA}}\\
  	\vspace{25mm}
  	{\large \textbf{ANTEPROYECTO}}\\
  	\vspace{15mm}
  	{Cuadro de Mando para la Visualizaci�n de la Calidad del Software, Integrado con un Entorno de Medici�n}\\
  	\vspace{15mm}
  	{Jose Domingo L�pez L�pez}
  	\vspace{20mm}
  	\begin{flushright}
  		\textbf{Septiembre, 2010}
  	\end{flushright}
	\end{center}
\end{titlepage}
% Derechos de autor
\mbox{}
\vfill{}
\begin{small}
	Cuadro de Mando para la Visualizaci�n de la Calidad del Software, Integrado con un Entorno de Medici�n
	
	\copyright~ 2010 Jose Domingo L�pez L�pez; Mois�s Rodr�guez Monje; Mar�a de los �ngeles Moraga de la Rubia
	
	En virtud del Art�culo 7 del Real Decreto Legislativo 1/1996, de 12 de abril por el que se aprueba el Texto Refundido de la Ley de Propiedad Intelectual, modificado por la Ley 23/2006 de 7 de julio, este PFC se considera una obra en colaboraci�n entre las diferentes partes. Por tanto la propiedad intelectual de este PFC de manera aislada, sin ser integrado con las herramientas y entornos de la empresa Alarcos Quality Center, ser� compartida con iguales porcentajes entre el alumno, el director y el tutor acad�mico.
	
	Este documento fue compuesto con \LaTeX{}. Im�genes generadas con GIMP 2.
\end{small}
\clearpage
% En los �ndices numeramos las p�ginas en n�meros romanos
\pagestyle{plain}
\pagenumbering{Roman}
% �ndice de secciones
\tableofcontents


\parskip=10pt
\pagenumbering{arabic}
\pagestyle{fancy}
% Aqui se incluyen los archivos .tex que forman el documento
\section{Introducci�n}

A�os atr�s, las organizaciones y empresas se preocupaban de la calidad de los procesos que se segu�an para desarrollar software bas�ndose en la premisa de que \textit{``con un proceso de calidad se obtiene un producto de calidad''} \cite{evaluacion-calidad}, pero dados los problemas que han surgido en sistemas ya implantados y en producci�n, y las p�rdidas que �stos han ocasionado, ha sido necesario actualizar la percepci�n que se tiene acerca de la calidad del software. Por ello, las organizaciones comenzaron a preocuparse por la calidad del producto software desde el punto de vista de los clientes y usuarios finales, abarcando un mayor n�mero de \textit{stakeholders}, con el fin de aumentar el grado de satisfacci�n y confianza por parte de los mismos.

Es por esto por lo que se tratado de asegurar la funcionalidad y usabilidad de los sistemas profundizando en �rea de \textit{testing}, pero se han quedado sin cubrir aspectos de seguridad y mantenibilidad. Y es que un producto software no es �nicamente el c�digo que se ejecuta sobre una m�quina, si no que lo son tambi�n los artefactos que se han ido produciendo a lo largo de todo el ciclo de vida: modelos, diagramas, documentos, etc.

Para paliar algunos de los problemas expuestos anteriormente, se comenzaron a utilizar dos est�ndares:
\begin{myitemize}
	\item \textbf{ISO/IEC 9126} \cite{iso9126}, que define el modelo de calidad de un producto, incluyendo las caracter�sticas, subcaracter�sticas y m�tricas (de calidad interna, calidad externa y calidad en uso) que han de tenerse en cuenta, pero sin proporcionar una metodolog�a para su evaluaci�n. Las caracter�sticas que propone son funcionalidad, fiabilidad, usabilidad, eficiencia, mantenibilidad y portabilidad.
	\item \textbf{ISO/IEC 14598} \cite{iso14598}, que proporciona gu�as y requisitos para el proceso de evaluaci�n, cubriendo as� los aspectos de los que carece la norma ISO/IEC 9126.
\end{myitemize}

Pero entre ambos est�ndares existen distintas incompatibilidades e inconsistencias -tales como diferencias en el vocabulario, t�rminos y conceptos, la necesidad de a�adir recomendaciones, metodolog�as, gu�as y nuevas necesidades para especificar las dimensiones de la calidad del software, entre otros- que motivan la creaci�n de la familia \textbf{ISO/IEC 25000}, denominada SQuaRe (\underline{S}oftware product \underline{Qua}lity \underline{R}equirements and \underline{E}valuation).

La norma ISO/IEC 25000 establece criterios para la especificaci�n de requisitos de calidad de productos software y su evaluaci�n, un modelo de calidad y m�tricas para los distintos atributos de dicho modelo. N�tese que no es necesario alcanzar una calidad perfecta o absoluta, sino la necesaria y suficiente para cada contexto de uso y en funci�n de las necesidades de los \textit{stakeholders}, que son todos las entidades involucradas en el producto: clientes, usuarios finales, factor�as y empresas de desarrollo, etc�tera.
% TODO En este p�rrafo sacar lo del post de mois�s
% http://www.iso15504.es/index.php/component/content/article/82-blog/165-la-familia-de-normas-iso-25000-y-su-relacion-con-iso-15504-spice.html

La aparici�n de esta norma, cuyo objetivo es reemplazar a las normas ISO/IEC 9126 e ISO/IEC 14598, junto con las nuevas necesidades que surgen en las empresas para controlar y evaluar la calidad de los desarrollos inform�ticos, motivan la aparici�n del proyecto \textbf{MEDUSAS} (\underline{M}ejora y \underline{E}valuaci�n del \underline{D}ise�o, \underline{U}sabilidad, \underline{S}eguridad y \underline{M}antenibilidad del \underline{S}oftware). MEDUSAS (Fig. \ref{fig:arch-medusas}) es un entorno metodol�gico e instrumental basado en la ISO/IEC 25000 para satisfacer dichas necesidades, y evaluar la calidad del software teniendo en cuenta todos y cada uno de sus artefactos: c�digo fuente, dise�o, diagramas, modelos, etc.

\imagen{./images/comp_medusas.png}{0.4}{Arquitectura del proyecto MEDUSAS}{fig:arch-medusas}

Aqu� es donde surge la motivaci�n del Proyecto Fin de Carrera (en adelante PFC), que consistir� en realizar un cuadro de mando para la visualizaci�n de la calidad del software, integrado con el entorno MEDUSAS. Como resultado de la integraci�n, se dispondr� de un sistema completo que permitir� automatizar la adquisici�n, presentaci�n y an�lisis \cite{gil95} de las m�tricas que definen la calidad del producto software.

\clearpage
\section{Objetivos}

Los objetivos que se esperan lograr con la realizaci�n de este PFC se han dividido en dos categor�as. En la primera de ellas, como objetivo principal, se introducir� el producto final que se espera obtener y, por �ltimo, como objetivos secundarios, se expondr�n los objetivos que se ir�n alcanzando en las distintas fases de elaboraci�n del PFC.


\subsection{Objetivo principal} \label{sec:objetivo-principal}

El PFC consiste en la elaboraci�n de una \textbf{aplicaci�n web} para la visualizaci�n de los datos obtenidos como resultado de las operaciones de medici�n y evaluaci�n de la calidad del producto software del entorno MEDUSAS, que est� basado en la nueva familia de normas ISO/IEC 25000. Para ello, se desarrollar� un proyecto Java que ser� capaz de representar gr�ficamente los distintos niveles de detalle procedentes de los resultados de evaluaci�n de la calidad del producto software.

Dicha aplicaci�n web ser� desarrollada utilizando las �ltimas tecnolog�as empleadas en desarrollo y dise�o web (Ver Secci�n \ref{sec:medios}) con el fin de obtener una RIA (Rich Internet Application) \cite{def:ria}.

Una RIA se fundamenta en una arquitectura cliente/servidor as�ncrona, segura y escalable, disminuyendo la carga de trabajo en el servidor y siendo accesible desde cualquier navegador web. As�, estar� disponible para cualquier persona desde cualquier rinc�n del planeta, tratando de maximizar la experiencia de usuario y producir un incremento en la productividad.


\subsection{Objetivos secundarios} \label{sec:objetivos-secundarios}

A lo largo del proceso de investigaci�n y elaboraci�n del PFC se alcanzar�n las siguientes metas:
\begin{myitemize}
	\item Comprender las actuales normas sobre la calidad del software.
	\item Investigar sobre las principales formas de visualizaci�n de calidad del software.
	\item Estudiar los principales entornos y medios actuales para la visualizaci�n de la calidad del software, as� como las librer�as gr�ficas que permiten su implementaci�n.
	\item Implementar un entorno que permita visualizar la calidad de un producto software de acuerdo a las caracter�sticas de calidad impuestas por la ISO/IEC 25000.
	\item Integrar dicho entorno de visualizaci�n con un entorno real de medici�n de la calidad del producto software.
\end{myitemize}
\clearpage
\section{M�todo y fases de trabajo}

Dada la naturaleza del proyecto, en el que se ir�n definiendo requisitos y funcionalidades a lo largo de todo su ciclo de vida, se ha optado por utilizar una metodolog�a de desarrollo de software gen�rica que permita adaptarse a este caso de estudio. Por estas razones, se ha seleccionado el \textbf{Proceso Unificado de Desarrollo} (en adelante PUD) como metodolog�a de trabajo.

El PUD \cite{jacobson2000proceso} es una evoluci�n del Proceso Unificado de Rational, que define un ``conjunto de actividades necesarias para transformar los requisitos de usuario en un sistema software''. Es un marco gen�rico que puede especializarse para una gran variedad de sistemas de software cualesquiera que sean el �rea de aplicaci�n o el tama�o del proyecto. Sus principales caracter�sticas son las siguientes:
\begin{myitemize}
	\item \textbf{Dirigido por casos de uso}. Para poder desarrollar un sistema es necesario saber qu� necesitan sus usuarios. Un usuario puede ser un ser humano u otro sistema que interacciona con el sistema que se est� desarrollando. Las necesidades de un usuario se denominan requisitos funcionales y se representan por medio de casos de uso.
	
	Los casos de uso gu�an el proceso de desarrollo desde la especificaci�n de requisitos hasta las pruebas y se utilizan para crear los modelos que permitan la construcci�n e implementaci�n de los mismos. Todos los casos de uso juntos constituyen el \textbf{modelo de casos de uso}.
	
	\item \textbf{Centrado en la arquitectura}. Un sistema software puede contemplarse desde varios puntos de vista. La arquitectura software incluye los aspectos est�ticos y din�micos m�s significativos del sistema y debe estar profundamente relacionada con los casos de uso ya que debe permitir el desarrollo de los mismos. Por esta raz�n, la arquitectura y los casos de uso deben desarrollarse en paralelo.
	
	\item \textbf{Iterativo e incremental}. Partiendo del dicho latino \textit{divide et impera} (en castellano, \textit{divide y vencer�s}), el ciclo de vida del proyecto (Fig. \ref{fig:mapa-conceptual-pud}) se divide en ciclos. Cada ciclo consta de cuatro fases que a su vez se dividen en iteraciones. Las fases de un ciclo son las siguientes:
	\begin{myitemize}
		\item \textbf{Inicio o Factibilidad}. Se obtiene un modelo de casos de uso simplificado, se identifican riesgos potenciales y se estima el proyecto de manera aproximada.
		
		\item \textbf{Elaboraci�n}. Se mejora el modelo de casos de uso y se dise�a la arquitectura del sistema. A continuaci�n, se desarrollan los casos de uso m�s cr�ticos que se identificaron en la fase de inicio y se obtiene una \textit{l�nea base} de la arquitectura para planificar las actividades y estimar los recursos necesarios para terminar el proyecto.
		
		\item \textbf{Construcci�n}. Se crea el producto en base a las dos fases anteriores.
		
		\item \textbf{Transici�n}. Implica la correcci�n de errores y \textit{bugs}, as� como el mantenimiento del sistema.
	\end{myitemize}
	
	 En cierto modo, una iteraci�n se corresponde con un mini-proyecto que resulta en un \textit{incremento}. En cada iteraci�n se aborda un conjunto de casos de uso que ampl�an la funcionalidad global del producto desarrollado y trata los riesgos m�s importantes. Si la iteraci�n cumple sus objetivos se contin�a con la pr�xima, en caso contrario se vuelve a las primeras etapas de la misma y se toman otras decisiones que permitan cumplir sus objetivos. Las iteraciones deben ser seleccionadas y ejecutarse de forma controlada y siguiendo un flujo de trabajo. Este flujo de trabajo consta de cinco etapas:
	\begin{myitemize}
		\item \textbf{Requisitos}. El equipo de desarrollo y los clientes establecen comunicaciones por diferentes medios para especificar qu� es lo que el usuario espera del sistema que se va a desarrollar.
		
		\item \textbf{An�lisis}. Se identifican y especifican los casos de uso relevantes.
		
		\item \textbf{Dise�o}. Se crea un dise�o utilizando la arquitectura seleccionada como gu�a para tratar de dotar al sistema de las funcionalidades representadas por los casos de uso identificados en la etapa de an�lisis.
		
		\item \textbf{Implementaci�n}. Se implementan las decisiones tomadas en la etapa de dise�o, con el fin de construir las funcionalidades que satisfacen los casos de uso identificados en la etapa de an�lisis.
		
		\item \textbf{Pruebas}. Se verifica que los cambios realizados en el sistema satisfacen los casos de uso.
	\end{myitemize}
\end{myitemize}


\imagenBorde{./images/mapa-conceptual-pud.png}{0.39}{Mapa conceptual del PUD}{fig:mapa-conceptual-pud}

\clearpage
\section{Medios que se pretenden utilizar} \label{sec:medios}

Los medios y recursos necesarios para el desarrollo y despliegue del PFC se categorizan en \textit{medios hardware} y \textit{medios software}.


\subsection{Medios hardware}

Tanto para el desarrollo como el despliegue del sistema bastar� con un ordenador sobremesa o port�til. Dicho ordenador deber� estar dotado de una conexi�n a Internet y actuar como servidor web para permitir que los usuarios puedan acceder a la aplicaci�n por medio de un navegador web convencional.

La tarjeta gr�fica del sistema requiere especial menci�n. �sta no tiene por qu� poseer una potencia extraordinaria, pero es de vital importancia que disponga de sus controladores gr�ficos debidamente instalados para una correcta ejecuci�n del motor gr�fico.


\subsection{Medios software}

Tanto la aplicaci�n web como el motor gr�fico ser�n desarrollados en tecnolog�as basadas en Java y, debido a su naturaleza \textit{multiplataforma}, se espera que no existan restricciones en cuanto al sistema operativo que el ordenador debe ejecutar. Sin embargo, �ste debe tener debidamente instalado el entorno de ejecuci�n Java (JRE o \textit{Java Runtime Environment})

Por otro lado, son necesarias las siguientes herramientas:
\begin{myitemize}
	\item JDK 1.4+ (\textit{Java Development Kit}).
	\item Un entorno integrado de desarrollo (IDE o \textit{Integrated Development Environment}). Se utilizar� \textbf{Eclipse EE Helios}.
	\item Un navegador web con JavaScript habilitado para que se pueda establecer una comunicaci�n entre la aplicaci�n web y el motor gr�fico, que estar� incrustado en un applet, y poder beneficiarse de las caracter�sticas de la tecnolog�a AJAX (\textit{Asynchronous JavaScript And XML}).
	\item Un sistema para la gesti�n y construcci�n de proyectos. Se har� uso de \textbf{Maven 2}.
	\item Un servidor web capaz de desplegar aplicaciones web con estructura WAR (\textit{Web-Archive} ). \textbf{Apache Tomcat 7} es la soluci�n m�s factible gracias a su licencia, eficiencia y rendimiento.
	\item Un sistema gestor de bases de datos. Se emplear� \textbf{MySQL Community Server 5.x}.
	\item \textbf{Visual Paradigm} para modelado y dise�o de diagramas UML.
	\item \LaTeX{} para la elaboraci�n de documentaci�n t�cnica asociada al proyecto.
	\item GIMP e Inkscape para el dise�o y edici�n de figuras, im�genes y gr�ficos.
\end{myitemize}

\clearpage

\pagestyle{plain}
% A�ade los ap�ndices al �ndice
\addcontentsline{toc}{section}{Anexos}
% Ap�ndices
\section*{Anexos}

\appendix

\section{Solicitud de evaluaci�n}

\newsavebox{\caja}
\sbox{\caja}{\textbf{Escuela Superior de Inform�tica}}
\newlength{\anchoCaja}
\setlength{\anchoCaja}{\wd\caja}
\noindent
\begin{minipage}{\anchoCaja}
	\begin{center}
	\noindent\includegraphics[keepaspectratio,scale=0.25]{./images/esi_bw.png}\\
		\noindent\textbf{Escuela Superior de Inform�tica}
	\end{center}
\end{minipage}

\vspace{10mm}
\noindent D. Jose Domingo L�pez L�pez, con D.N.I. n\textordmasculine{} 71219116-F, alumno de la Escuela Superior de Inform�tica.

\vspace{5mm}
\noindent EXPONE:\\
\textemdash{} Que se encuentra matriculado del PFC.

\vspace{5mm}
\noindent SOLICITA:\\
\textemdash{} Que la Comisi�n Acad�mica eval�e el anteproyecto de fin de carrera que acompa�a, titulado \textit{Cuadro de Mando para la Visualizaci�n de la Calidad del Software, Integrado con un Entorno de Medici�n}.

\vspace{10mm}
\noindent Ciudad Real, a \today{}\\

\newlength{\factor}
\setlength{\factor}{0.32\textwidth}
\noindent 
\begin{minipage}{\factor}
El alumno,\\
\\
\\
\\
Fdo.:
\end{minipage}
\hfill
\begin{minipage}{\factor}
V\textordmasculine B\textordmasculine{} El director,\\
\\
\\
\\
Fdo.:
\end{minipage}
\hfill
\begin{minipage}{\factor}
V\textordmasculine B\textordmasculine{} El tutor acad�mico\\
\\
\\
\\
Fdo.:
\end{minipage}

\vspace{10mm}
\begin{center}
	\noindent SR. DIRECTOR DE LA ESCUELA SUPERIOR DE INFORM�TICA
\end{center}

\clearpage

\section{Contrato de Propiedad Intelectual}\label{anexo:cpi}

\begin{flushright}
	En Ciudad Real, a \today	
\end{flushright}

\begin{center}
	\textbf{REUNIDOS}
\end{center}

DE UNA PARTE, \textbf{D. Jose Domingo L�pez L�pez} mayor de edad, con D.N.I. n�mero \textbf{71219116-F}, Alumno de la Escuela Superior de Inform�tica, y en adelante el \textbf{``ALUMNO DEL PFC''}.

DE OTRA PARTE, \textbf{D�a. Mar�a de los �ngeles Moraga de la Rubia} mayor de edad, con D.N.I. n�mero \textbf{5683662-V}, Profesora de la Escuela Superior de Inform�tica, y en adelante el \textbf{``TUTOR ACAD�MICO DEL PFC''}.

DE OTRA PARTE, \textbf{D. Mois�s Rodriguez Monje} mayor de edad, con D.N.I. n�mero \textbf{5928188-F}, Director T�cnico de la empresa Alarcos Quality Center (en adelante AQC), y en adelante el \textbf{``DIRECTOR DEL PFC''}.

Todos ellos reconoci�ndose mutuamente capacidad suficiente para la celebraci�n del presente Contrato,

\begin{center}
	\textbf{EXPONEN}	
\end{center}

\noindent\textbf{PRIMERO:} El objeto del presente contrato es la determinaci�n de la propiedad intelectual del PFC titulado \textbf{``Cuadro de Mano para la Visualizaci�n de la Calidad del Software, Integrado con un Entorno de Medici�n''}, y del resultado de la integraci�n del mismo con las herramientas usadas por \textbf{AQC}. As� como, establecer el compromiso de la confidencialidad que el \textbf{ALUMNO DEL PFC} contraer� con \textbf{AQC}.

\noindent\textbf{SEGUNDO:} La propiedad intelectual del producto software resultante del \textbf{PFC} integrado con las herramientas y entornos de la empresa \textbf{AQC} pertenecer� �nica y exclusivamente a \textbf{AQC}, no puediendo el \textbf{ALUMNO DEL PFC} utilizarlas sin el consentimiento expreso de la empresa. Sin embargo, la propiedad intelectual del \textbf{PFC} de manera aislada, as� como de los documentos, diagramas, modelos, etc. que lo acompa�en ser� compartida a partes iguales entre el \textbf{ALUMNO DEL PFC}, el \textbf{TUTOR DEL PFC} y el \textbf{DIRECTOR DEL PFC}.

\noindent\textbf{TERCERO:} El \textbf{ALUMNO DEL PFC} trabajar� para la realizaci�n de su \textbf{PFC} con documentaci�n, productos software y entornos de la empresa \textbf{AQC} y por ello, por medio de la presente, se compromete de forma irrevocable ante \textbf{AQC} a:
\begin{enumerate}
	\item Asumir las pol�ticas, los objetivos y las directrices establecidas por \textbf{AQC}, y particularmente la pol�tica de calidad, las directrices y procedimientos incluidos o referenciados en el manual de calidad de \textbf{AQC}.
	
	\item Actuar de forma que se preserve una imagen positiva y eficad de \textbf{AQC} guardando al debida diligencia en relaci�n con todas las labores efectuadas, actuando en todo momento de buena fe y en aras al leal cumplimiento de mis obligaciones para con \textbf{AQC}.
	
	\item Evitar intervenir en cualquier actividad que pueda disminuir la confianza en la competencia, imparcialidad, juicio o integridad operativa de \textbf{AQC}.
	
	\item Asegurar la imparcialidad y estar libre de toda presi�n indebida, comercial, financiera o de otra �ndole, que pueda influir en el juicio t�cnico empleado en la realizaci�n de sus actividades en \textbf{AQC}.
	
	\item No alegar su calidad de colaborador de \textbf{AQC} con fines ajenos a las misiones que se le hayan confiado por \textbf{AQC}. No admitir ning�n tipo de retribuci�n ajena a \textbf{AQC} por cualquier actividad realizada en nombre de \textbf{AQC} o en calidad de persona cualificada por \textbf{AQC}.
	
	\item En particular, durante la realizaci�n de las actividades del \textbf{PFC}:
	\begin{itemize}
		\item Detectar desviaciones / no conformidades basadas en evidencias objetivas y no en opiniones pesonales o valoraciones subjetivas.
		
		\item No adoptar nunca actitudes de prepotencia ante la empresa.
		
		\item No tener inconveniente en modificar una opini�n si la empresa le demuestra que estaba equivocado.
		
		\item No ceder ante presiones o coacciones que pueda recibir si est� convencido de sus criterios y no le demuestran con contrario.
		\end{itemize}
		
	\item Participar, en la medida de los posible, en las actividades, reuniones o cursos de armonizaci�n de criterios siempre y cuando sea requerido por \textbf{AQC}. Compartir con el personal de \textbf{AQC}, su conocimiento, experiencia, ideas y sugerencias.
	
	\item Avisar a \textbf{AQC} de cualquier cambio o interrupci�n en sus actividades profesionales.
	
	\item Durante su relaci�n con \textbf{AQC} y tras ella, guardar el m�s estricto secreto sobre todos los asuntos, documentos y registros confidenciales de car�cter t�cnico, organizativo, comercial y econ�mico de \textbf{AQC} y aquella que pudiera tener conocimiento dentro de marco de sus actividades. Tratando como documentaci�n confidencial:
	\begin{itemize}
		\item Solicitud, formal o informal, de un proyecto. Queda incluido el hecho de la existencia de la solicitud.
		
		\item La entrega por una organizaci�n que ha solicitado o ha realizado un proyecto con \textbf{AQC}.
		
		\item Los datos de los productos o servicios o de las organizaciones que ha solicitado o ha realizado un proyecto con \textbf{AQC}.
		
		\item Toda la correspondencia entre \textbf{AQC} y las organizaciones con que trabaja o tiene relaci�n.
		
		\item Toda la correspondencia entre \textbf{AQC} y sus subcontratistas y viceversa.
		
		\item La que los documentos del sistema se�alen como confidencial.
		
		\item Los documentos del sistema de gesti�n de \textbf{AQC}.
		
		\item Todos los documentos que \textbf{AQC} le facilite para la realizaci�n exitosa de su \textbf{PFC}, as� como las herramientas y modelos con los que el \textbf{ALUMNO DEL PFC} deba integrar su trabajo.
	\end{itemize}
	
	\item Adoptar todas las precauciones para evitar que se divulguen, directa o indirectamente, por causa propia o de las personas de las que es responsable, documentos o informaciones de los que pudiera tener conocimiento dentro del marco de sus actividades en los proyectos de \textbf{AQC}. Estas precauciones incluyen:
	\begin{itemize}
		\item Manipular y mantener debidamente archivada la documentaci�n.
		\item No dejar incontrolada documentaci�n confidencial en salas de reuniones y visitas.
		\item Evitar que personas ajenas accedan a lugares de trabajo y archivos.
	\end{itemize}
	
	\item En caso de acceso y tratamiento de datos de car�cter personal, por exigencias el trabajo que desempe�e, asegurar� la confidencialidad de los mismos respetando y cumpliendo en todo momento los procedimientos y directrices instaurados en \textbf{AQC} al efecto y la legislaci�n vigente aplicable.
	
	\item Respetar la titularidad de \textbf{AQC} sobre cualquier proyecto, informe, documento e im�genes que realice en el desempe�o de las tareas que se ele encomiendan durante el \textbf{PFC} (respetando y protegiendo los derechos de propiedad intelectual o industrial de \textbf{AQC}).
\end{enumerate}

Y en prueba de cuando antecede, las Partes suscriben el Contrato, en cuatro ejemplares y a un solo efecto, en el lugar y fecha se�alados en el encabezamiento.

\noindent
\begin{minipage}[t][6cm][c]{1.0\textwidth}
	\phantom{EL ALUMNO DEL PFC}
	\hfill
	\phantom{EL TUTOR DEL PFC}
	\hfill
	\phantom{EL DIRECTOR DEL PFC} \\
	\textbf{EL ALUMNO DEL PFC}
	\hfill
	\textbf{EL TUTOR DEL PFC}
	\hfill
	\textbf{EL DIRECTOR DEL PFC}
\end{minipage}


\clearpage\phantomsection
% A�ade la bibliograf�a al �ndice
\addcontentsline{toc}{section}{\bibname}
% Bibliograf�a
\bibliographystyle{plain}
\pagestyle{plain} 
\bibliography{Bibliografia}

\end{document}
